\documentclass[11pt]{article}
\usepackage[T1]{fontenc}
\usepackage{lmodern,amsmath,amssymb,amsthm,mathtools,microtype,booktabs}
\usepackage[letterpaper,margin=1.05in]{geometry}
\usepackage{xcolor}
\usepackage{hyperref}
\hypersetup{colorlinks=true,linkcolor=blue!50!black,urlcolor=blue!50!black,citecolor=blue!50!black}
\newtheorem{theorem}{Theorem}[section]
\newtheorem{lemma}[theorem]{Lemma}
\newtheorem{corollary}[theorem]{Corollary}
\newtheorem{proposition}[theorem]{Proposition}
\theoremstyle{remark}\newtheorem{remark}[theorem]{Remark}
\newcommand{\Z}{\mathbb Z}\newcommand{\R}{\mathbb R}\newcommand{\F}{\mathbb F}
\newcommand{\covol}{\operatorname{covol}}
\newcommand{\diag}{\operatorname{diag}}
\newcommand{\tr}{\operatorname{tr}}
\newcommand{\norm}[1]{\left\lVert #1\right\rVert}
\newcommand{\bigo}{\mathcal O}
\allowdisplaybreaks[2]
\setlength{\emergencystretch}{2em}
\title{\textbf{Sharp growth of integral Gram determinants and the generic modular column-number problem}}
\author{Research note for ProveIt}
\date{27 November 2026}
\begin{document}
\maketitle
\begin{abstract}
For an integer matrix $A$ of size $r\times n$ with all maximal minors nonzero, let $D(A)$ be the gcd of the maximal minors and let $E_r(n)$ minimize $\det(AA^{\mathsf T})/D(A)^2$. We prove $E_r(n)=\Theta_r(n^{3r-2})$ for every fixed $r\ge2$, answering the order-of-growth part of the question \texttt{src:q:gram} in ProveIt's article on Dirichlet GRH. The lower bound is a consequence of a homogeneous hyperplane cover for anisotropic integer boxes, combined with a short dual lattice basis and an isotropic change of metric. For the upper bound, short integer lifts of the projective moment curve over a prime field suffice. The same method gives an affirmative answer to Question 3.4 of Kriepke, Kyureghyan and Schymura: the largest number $g(\Delta,r)$ of columns of a generic integer matrix with largest absolute maximal minor \emph{exactly} $\Delta$ has order $\Theta_r(\Delta^{1/(r-1)})$. We handle the exact-value condition for every sufficiently large integer $\Delta$ using a prescribed determinant, not just an infinite subsequence. The results are unconditional; the transfer concerns the Gram/covolume mechanism, not the analytic zero-location theorem. Exact finite certificates accompany the proofs.
\end{abstract}

\section{Questions and statements}
A matrix $A=[v_1\ \cdots\ v_n]\in\Z^{r\times n}$, $n\ge r\ge2$, is called \emph{generic} here if every $r$ distinct columns are independent. Write $A_I$ for its submatrix on a set $I$ of $r$ column indices, and put
\begin{equation}\label{eq:def}
 D(A)=\gcd_{|I|=r}|\det A_I|,\qquad
 \mathcal E(A)=\frac{\det(AA^{\mathsf T})}{D(A)^2},\qquad
 E_r(n)=\min_{A\ \mathrm{generic}}\mathcal E(A).
\end{equation}
The quotient is a positive integer, by Cauchy--Binet:
\begin{equation}\label{eq:CB}
 \mathcal E(A)=\sum_{|I|=r}\left(\frac{\det A_I}{D(A)}\right)^2.
\end{equation}
There are generic integer matrices of every such size (for instance ordinary Vandermonde matrices), so the minimum exists. The merged ProveIt article \cite{proveit}, in \texttt{src:cor:gramrank}, obtains a lower bound of order $n^{r+2}$ at fixed $r$ by contracting to rank two. It explicitly asks in \texttt{src:q:gram} for the order of $E_r(n)/\binom nr$ at $r\ge3$. Its rank-two analysis gives a much more precise constant in that special dimension, which we do not seek to replace.

\begin{theorem}[Gram energy]\label{thm:gram}
For each fixed integer $r\ge2$, there is an explicit $c_r>0$ such that for every $n\ge2r$,
\begin{equation}\label{eq:bounds}
 c_r n^{3r-2}\le E_r(n)\le 2^{2r-2}n^{3r-2}.
\end{equation}
In particular
\[
 \boxed{\displaystyle\frac{E_r(n)}{\binom nr}=\Theta_r(n^{2r-2}).}
\]
The lower bound applies to each generic matrix, including arbitrarily anisotropic ones. A stronger version in terms of hyperplane occupancy is given in Theorem~\ref{thm:robust}.
\end{theorem}

The second target is a published question. For $\Delta\in\Z_{>0}$, let
\begin{align*}
 g(\Delta,r)&=\max\{n:\ A\in\Z^{r\times n}\text{ generic},\ \max_{|I|=r}|\det A_I|=\Delta\},\\
 g_{\le}(\Delta,r)&=\max\{n:\ A\in\Z^{r\times n}\text{ generic},\ \max_{|I|=r}|\det A_I|\le\Delta\}.
\end{align*}
We distinguish exact and bounded versions because they are not logically interchangeable. Kriepke, Kyureghyan and Schymura \cite[Theorem 1.4, Proposition 3.3 and Question 3.4]{kks} give upper exponent $2/r$, construct matrices demonstrating lower exponent $1/(r-1)$, and ask whether the lower exponent is the correct one. Here is the affirmative answer, with a uniform-in-$\Delta$ construction for the exact version.

\begin{theorem}[Generic column number]\label{thm:column}
For every fixed $r\ge2$,
\begin{equation}\label{eq:g}
 \boxed{g(\Delta,r)=\Theta_r(\Delta^{1/(r-1)})},\qquad
 g_{\le}(\Delta,r)=\Theta_r(\Delta^{1/(r-1)})
\end{equation}
as the integer $\Delta\to\infty$. Moreover, for every sufficiently large $\Delta$, there is a matrix realizing the lower order with precisely one maximal minor of absolute value $\Delta$; all other maximal minors are nonzero and have absolute value at most $\Delta/2$.
\end{theorem}
For rank three, for instance, the exponent decreases from the previously published upper exponent $2/3$ to $1/2$. This is a fixed-rank result, not a bound uniform in growing $r$ and not a claim about matrices with vanishing maximal minors. It does not settle the separate conjecture on polynomial-time solvability of general modular integer programs.

\section{A homogeneous hyperplane cover}
The key estimate is a version of a short integer-relation argument that does not assume the box is roughly cubic. It is convenient to state the result independently of matrices. A homogeneous hyperplane is the kernel of a nonzero real linear functional. Here all the covering functionals can be taken integral.

\begin{lemma}[Anisotropic box covering]\label{lem:box}
Let $L_i\ge1$ be real numbers, $1\le i\le r$, $r\ge2$, and put $P=\prod_i L_i$. The integer points in the box $\prod_i[-L_i,L_i]$ are covered by at most
\begin{equation}\label{eq:cover}
 K_r P^{1/(r-1)},\qquad K_r=(2r+3)^r,
\end{equation}
homogeneous hyperplanes with nonzero integer normals. In particular, if a family of indexed integer points in that box has at most $h$ members on each homogeneous hyperplane, it has at most $h K_r P^{1/(r-1)}$ members.
\end{lemma}

\begin{proof}
Permute the coordinates so that $L_1\le\cdots\le L_r$. For $2\le j\le r$, define
\[
 R_j=(L_1\cdots L_j)^{1/(j-1)}.
\]
Choose $j$ minimizing $R_j$, and write $R=R_j$. Then $1\le R\le P^{1/(r-1)}$ and $R\ge L_j$. Indeed $R_2=L_1L_2\ge L_2$. For $j>2$, if $L_j>R_j$ then
\[
 R_{j-1}^{j-2}=\frac{R_j^{j-1}}{L_j}<R_j^{j-2},
\]
contradicting minimality. Thus $R/L_i\ge1$ on the selected coordinates.

Set $C=r+1$ and $m_i=\lfloor CR/L_i\rfloor$ for $1\le i\le j$. Consider the fixed set of coefficient vectors
\[
 \mathcal B=\prod_{i=1}^j\{0,1,\ldots,m_i\}.
\]
For every integer point $z$ in the original box, the integers $a\cdot(z_1,\ldots,z_j)$, $a\in\mathcal B$, range over an interval of length at most $\sum_i m_i|z_i|\le jCR$. There are at most $jCR+1$ possible values. On the other hand,
\[
 |\mathcal B|=\prod_i(m_i+1)>\frac{(CR)^j}{L_1\cdots L_j}=C^jR>jCR+1.
\]
The last inequality follows from $R\ge1$, $j\ge2$ and $C=r+1$: $C^j-jC\ge C(C-j)\ge3$. The pigeonhole principle gives distinct $a,a'$ with equal dot product. Their difference, padded with zero entries in the last $r-j$ coordinates, is a nonzero integer normal $u$ with $u\cdot z=0$ and $|u_i|\le m_i$.

The set of possible normals is independent of $z$. Its cardinality, even including zero and both signs, is at most
\[
 \prod_{i=1}^j(2m_i+1)\le(2C+1)^j\frac{R^j}{L_1\cdots L_j}
 =(2C+1)^jR\le K_r P^{1/(r-1)}.
\]
The cover is established; summing the occupancy bound over its hyperplanes proves the last assertion.
\end{proof}

\begin{remark}
For a cube of halfwidth $H$ the cardinality is $O_r(H^{r/(r-1)})$. The subtlety is the choice of $j$. In a very thin box a normal involving all coordinates need not be available at the natural volume scale; one must allow a normal supported on a smaller subset. This is the reason the result works after an arbitrary change of lattice metric.
\end{remark}

\section{From boxes to Gram determinants}
We record a short-basis result with an elementary proof so that there is no dependence on a deep lattice-reduction theorem. The lattice covolume is the absolute determinant of a basis, computed in its Euclidean span.

\begin{lemma}[A short lattice basis]\label{lem:basis}
For every full-rank Euclidean lattice $\Lambda\subset\R^r$ there is a lattice basis $u_1,\ldots,u_r$ such that
\begin{equation}\label{eq:short}
 \prod_{i=1}^r\norm{u_i}_2\le M_r\covol(\Lambda),\qquad
 M_r=(2/\sqrt3)^{r(r-1)/2}.
\end{equation}
\end{lemma}
\begin{proof}
The dimension-one assertion is equality. Choose a shortest nonzero vector $u$ in $\Lambda$. It is primitive, since a nontrivial integer multiple would not be shortest. A primitive lattice vector extends to an integral lattice basis (in any existing basis its coordinate gcd is one, and elementary integer row operations reduce it to a standard basis vector). Orthogonal projection onto $u^\perp$ therefore gives a full-rank lattice $\Lambda'$ in dimension $r-1$, and
\[
 \covol(\Lambda)=\norm u_2\,\covol(\Lambda').
\]
Every nonzero $w\in\Lambda'$ has a lift $v$ whose component along $u$ has length at most $\norm u_2/2$, by subtracting a suitable integer multiple of $u$. Shortness of $u$ implies
\[
 \norm w_2^2+\norm u_2^2/4\ge\norm v_2^2\ge\norm u_2^2,
\]
so $\norm w_2\ge(\sqrt3/2)\norm u_2$. In particular such a lift satisfies $\norm v_2\le(2/\sqrt3)\norm w_2$. Apply the induction hypothesis to a basis of $\Lambda'$ and use adjusted lifts of its members together with $u$; these form a basis of $\Lambda$, because the projection kernel is $\Z u$. Multiplying the norm bounds gives the constant $M_r=(2/\sqrt3)^{r-1}M_{r-1}$.
\end{proof}

\begin{lemma}[Normalization]\label{lem:norm}
For any rank-$r$ integer matrix $A$, let $L=\sum_i\Z v_i$. A basis matrix $B$ for $L$ satisfies $|\det B|=D(A)$, where zero minors may be omitted from the gcd. If $A=BC$, then $C$ is integer, its columns generate $\Z^r$, $D(C)=1$, and
\[
 \det(CC^{\mathsf T})=\mathcal E(A).
\]
Homogeneous hyperplane occupancies and vanishing patterns of minors are unchanged.
\end{lemma}
\begin{proof}
All columns of $A$ are in $L$, so $C$ is integer and its columns generate $\Z^r$. Consequently there is an integer $n\times r$ matrix $J$ with $CJ=I_r$. Cauchy--Binet applied to $CJ$ shows that the gcd of the maximal minors of $C$ divides one. Each maximal minor of $A$ is $\det B$ times the corresponding minor of $C$. The gcd statement follows, and so does the determinant formula by taking determinants of $AA^{\mathsf T}=BCC^{\mathsf T}B^{\mathsf T}$. The transformation is invertible over $\R$.
\end{proof}

\begin{theorem}[Gram bound with hyperplane occupancy]\label{thm:robust}
Let $A\in\Z^{r\times n}$ have rank $r\ge2$. Let $h$ be the maximum number of its indexed columns in a homogeneous hyperplane of $\R^r$. If $h\le n/2$, then
\begin{equation}\label{eq:robust}
 \mathcal E(A)\ \ge\ \frac{n^r}{M_r^2(2r)^r}\left(\frac{n}{2hK_r}\right)^{2r-2}.
\end{equation}
\end{theorem}
\begin{proof}
Apply Lemma~\ref{lem:norm}, and denote the resulting normalized matrix again by $A$. Let $G=AA^{\mathsf T}$ and $w_i=G^{-1/2}v_i$. Positivity of $G$ follows from the full row rank. The isotropy identity is
\[
 \sum_i w_iw_i^{\mathsf T}=I_r,\qquad \sum_i\norm{w_i}_2^2=r.
\]
With $R=\sqrt{2r/n}$, strictly more than $n/2$ columns satisfy $\norm{w_i}_2\le R$ (if at least $n/2$ were strictly longer, their squared norms would already sum to more than $r$). Let $I$ be this set. Since $|I|>h$, these columns span $\R^r$.

The $w_i$ lie in the lattice $\Lambda=G^{-1/2}\Z^r$. Its dual, defined by integral pairing with $\Lambda$, is $\Lambda^*=G^{1/2}\Z^r$, and has covolume $\sqrt{\det G}$. By Lemma~\ref{lem:basis} choose a basis $u_1,\ldots,u_r$ of this dual lattice with product of lengths at most $M_r\sqrt{\det G}$. For $i\in I$, form the integer coordinate vector
\[
 z_i=(\langle u_1,w_i\rangle,\ldots,\langle u_r,w_i\rangle)^{\mathsf T}.
\]
These vectors belong to the box of halfwidths $L_j=R\norm{u_j}_2$. Each $L_j\ge1$: otherwise every corresponding integer coordinate is zero, contradicting spanning of the $w_i$ in $I$. The product satisfies
\begin{equation}\label{eq:product}
 P=\prod_j L_j\le (2r/n)^{r/2}M_r\sqrt{\det G}.
\end{equation}
The map from $w$ to $z$ is invertible, so each homogeneous hyperplane contains at most $h$ indexed members of this new family. Lemma~\ref{lem:box} now gives
\[
 n/2< |I|\le hK_rP^{1/(r-1)}.
\]
Combining with \eqref{eq:product} and squaring yields \eqref{eq:robust}, since $\det G=\mathcal E(A)$ by normalization.
\end{proof}

A generic matrix has $h=r-1$: no $r$ columns can lie in a hyperplane, whereas any $r-1$ columns do lie in one. Thus Theorem~\ref{thm:robust} gives the lower bound of Theorem~\ref{thm:gram} for $n\ge2r$, with
\begin{equation}\label{eq:cr}
 c_r=\left[M_r^2(2r)^r\bigl(2(r-1)K_r\bigr)^{2r-2}\right]^{-1}.
\end{equation}
These constants are deliberately nonoptimal. The exponent is sharp.

\begin{corollary}[A bounded number of vanishing minors]\label{cor:zeros}
Fix $r\ge2$ and a nonnegative integer $z$. The minimum normalized Gram determinant over rank-$r$ integer matrices of size $r\times n$ with \emph{at most} $z$ vanishing maximal minors has order $\Theta_{r,z}(n^{3r-2})$. More generally the lower bound \eqref{eq:robust} applies with $h$ replaced by
\[
 H_z=r-1+(r!z)^{1/r}
\]
whenever $H_z\le n/2$.
\end{corollary}
\begin{proof}
A homogeneous hyperplane containing $h$ columns accounts for $\binom hr$ vanishing minors. For $h\ge r$ this number is at least $(h-r+1)^r/r!$, so $h\le H_z$; the same inequality is immediate for $h<r$. Apply Theorem~\ref{thm:robust}. The generic construction below supplies the upper bound for the class with at most $z$ vanishing minors. This is not a claim about the optimum with exactly $z$ vanishing minors, or about the sharp leading constant.
\end{proof}

\section{Short projective lifts and the matching upper bound}
The following lemma is elementary but crucial. It uses the freedom to rescale each column independently modulo a prime. Merely lifting an affine moment curve coordinate by coordinate would be too expensive for the Gram energy.

\begin{lemma}[Short projective lift]\label{lem:lift}
Let $p$ be prime, $r\ge2$, and $\ell\subset\F_p^r$ a one-dimensional subspace. There is a vector $v\in\Z^r$ with $v\bmod p\in\ell\setminus\{0\}$ and
\[
 \norm v_\infty\le H=\lfloor p^{(r-1)/r}\rfloor.
\]
\end{lemma}
\begin{proof}
The quotient $\F_p^r/\ell$ has $p^{r-1}$ elements. The grid $\{0,\ldots,H\}^r$ contains $(H+1)^r>p^{r-1}$ points, so two distinct grid points have the same image in that quotient. Their difference $v$ satisfies the norm bound and reduces to an element of $\ell$. Since $H<p$, a difference reducing to zero would be zero as an integer vector, impossible for two distinct grid points.
\end{proof}

\begin{proposition}[Full-spark short columns]\label{prop:curve}
For every prime $p\ge n\ge r\ge2$, there is a generic $r\times n$ integer matrix $A$ with entries bounded by $H=\lfloor p^{(r-1)/r}\rfloor$ in absolute value. It satisfies
\begin{equation}\label{eq:upper}
 \mathcal E(A)\le\det(AA^{\mathsf T})\le n^r H^{2r}\le n^r p^{2r-2}.
\end{equation}
\end{proposition}
\begin{proof}
Choose distinct parameters $t_1,\ldots,t_n\in\F_p$ and write $q(t)=(1,t,\ldots,t^{r-1})^{\mathsf T}$. Apply Lemma~\ref{lem:lift} to the line spanned by each $q(t_i)$. The resulting $v_i$ reduce to $k_iq(t_i)$ for some nonzero $k_i\in\F_p$. Every maximal minor is nonzero modulo $p$, by the Vandermonde determinant formula, and hence is nonzero over $\Z$.

Each column has squared length at most $rH^2$, so $\tr(AA^{\mathsf T})\le nrH^2$. The geometric mean of the $r$ positive eigenvalues is at most their arithmetic mean. Therefore $\det(AA^{\mathsf T})\le(nH^2)^r$. Division by the squared gcd cannot increase it. This proves \eqref{eq:upper}.
\end{proof}

Bertrand's postulate supplies a prime $n\le p<2n$ for every $n\ge2$. Proposition~\ref{prop:curve} and the lower bound in Section~3 complete the proof of Theorem~\ref{thm:gram}.

\section{Answer to the generic column-number question}
For the upper bound, let a generic $r\times n$ matrix have all maximal minors bounded in absolute value by $\Delta$. If $n\ge2r$, Cauchy--Binet and Theorem~\ref{thm:gram} imply
\[
 c_r n^{3r-2}\le\mathcal E(A)\le \binom nr\Delta^2\le n^r\Delta^2.
\]
Hence
\begin{equation}\label{eq:gupper}
 n\le c_r^{-1/(2r-2)}\Delta^{1/(r-1)}.
\end{equation}
Matrices with $n<2r$ are covered by enlarging the constant. This proves the required upper bound for both versions, and in particular the maximum defining each is finite. For the bounded version, Proposition~\ref{prop:curve} already produces matrices of the required size, since each maximal minor is at most $r^{r/2}p^{r-1}$ by Hadamard's inequality.

For the exact version an additional argument is necessary. The construction below makes one anchor minor equal to an arbitrary prescribed integer $\Delta$. First construct the anchor basis with entries of order $\Delta^{1/r}$. Then choose the projective moment curve in coordinates adapted to that basis modulo a suitable prime.

\begin{lemma}[Balanced prescribed determinant]\label{lem:anchor}
For every positive integer $\Delta$ and $r\ge2$ there is an integer square matrix $B$ of order $r$ with $\det B=\Delta$ and each column of Euclidean length at most $2\sqrt r\,\Delta^{1/r}$.
\end{lemma}
\begin{proof}
Put $b=\lfloor\Delta^{1/r}\rfloor+1$ so $\Delta<b^r$ and $b\le 2\Delta^{1/r}$. Write the base-$b$ expansion $\Delta=\sum_{i=0}^{r-1}d_i b^i$, with $0\le d_i<b$. Let the first $r-1$ columns of $B$ be $b e_j-e_{j+1}$ for $j=1,\ldots,r-1$; the last column is $(d_0,\ldots,d_{r-1})^{\mathsf T}$. Expanding along the last column gives $\det B=\sum_{i=0}^{r-1}d_i b^i=\Delta$. Every entry has absolute value at most $b$, giving the stated norm bound.
\end{proof}

\begin{lemma}[A prime avoiding the anchor determinant]\label{lem:prime}
Fix $r\ge2$ and $a>0$. For all sufficiently large integers $\Delta$, there is a prime $p\nmid\Delta$ with
\begin{equation}\label{eq:prime}
 2^{-r-1}a\Delta^{1/(r-1)}\le p\le a\Delta^{1/(r-1)},\qquad p\ge r.
\end{equation}
\end{lemma}
\begin{proof}
Let $X=a\Delta^{1/(r-1)}$ and $m=\lfloor X/2^r\rfloor$. For large $\Delta$, $m\ge\max(2,r)$ and $m\ge X/2^{r+1}$. Apply Bertrand's postulate to $m,2m,\ldots,2^{r-1}m$. This gives $r$ distinct primes in disjoint open intervals, all greater than $m$ and less than $2^rm\le X$. For sufficiently large $\Delta$, $m^r>\Delta$ because $r/(r-1)>1$. The product of these primes is then greater than $\Delta$, so they cannot all divide it. Choose one that does not.
\end{proof}

\begin{proposition}[Exact modularity for every large determinant]\label{prop:exact}
Fix $r\ge2$. For all sufficiently large integers $\Delta$, there is a generic integer matrix with $p\ge b_r\Delta^{1/(r-1)}$ columns, where $b_r=2^{-r-1}(4r)^{-r}$, whose only maximal minor of absolute value $\Delta$ consists of its first $r$ columns. Every other maximal minor has absolute value at most $\Delta/2$.
\end{proposition}
\begin{proof}
Use Lemma~\ref{lem:prime} with $a=(4r)^{-r}$ and take the prescribed matrix $B$ of Lemma~\ref{lem:anchor}. Its reduction modulo $p$ is invertible by $p\nmid\Delta$. The vectors $q(0),\ldots,q(r-1)$ of the projective moment curve are a basis of $\F_p^r$. Thus the unique invertible linear transformation $T$ sending these vectors to the reductions of the columns of $B$ is
\[
 T=\bar B[q(0)\ \cdots\ q(r-1)]^{-1}\in\mathrm{GL}_r(\F_p).
\]
Keep the actual $r$ anchor columns of $B$. For every $t\in\F_p\setminus\{0,\ldots,r-1\}$, use Lemma~\ref{lem:lift} for the line $\F_p Tq(t)$ to append a short integer column. The full list has $p$ columns. Modulo $p$ the columns are precisely the $Tq(t)$ multiplied by nonzero scalars (one at each anchor). Every maximal minor is consequently nonzero modulo $p$, even those using anchors.

It remains to bound the absolute integer determinants. Every anchor has Euclidean length at most $2\sqrt r\,\Delta^{1/r}$. Every short column has length at most
\begin{align*}
 \sqrt r\,p^{(r-1)/r}
 &\le \frac{\sqrt r}{(4r)^{r-1}}\Delta^{1/r}\\
 &\le \frac{1}{2(2\sqrt r)^{r-1}}\Delta^{1/r}.
\end{align*}
For the last inequality the ratio of its left constant to its right constant is $(2\sqrt r)^{2-r}\le1$. The short norm bound is smaller than the anchor norm bound. Every minor other than the anchor minor uses at least one short column. By Hadamard's inequality its absolute value is at most
\[
 (2\sqrt r\,\Delta^{1/r})^{r-1}\,\frac{\Delta^{1/r}}{2(2\sqrt r)^{r-1}}=\frac{\Delta}{2}.
\]
The anchor determinant is exactly $\Delta$, proving the assertion. The lower bound on $p$ is \eqref{eq:prime}.
\end{proof}

Together with \eqref{eq:gupper}, Proposition~\ref{prop:exact} completes Theorem~\ref{thm:column}. No assumption about monotonicity of $g(\Delta,r)$ is used. In particular it is not necessary to extrapolate a lower bound on a special sequence of exact determinant values.

\section{Finite certificates and asymptotic interpretation}
The proof is symbolic. The companion Python program provides exact-integer tests of its finite constructions, with no external dependencies. For each projective direction it scans the nonzero scalar multiples modulo $p$, takes centered residues, and minimizes the sup norm, breaking ties deterministically. It verifies each determinant by fraction-free elimination and separately against the modular Vandermonde identity. It also computes the gcd, Gram determinant, sum of squared minors, and largest absolute minor. These examples are upper witnesses, not assertions of optimality.

\begin{center}
\begin{tabular}{rrrrrr}
\toprule
$r$&$n$&$p$&$H$&$\mathcal E(A)$&$\max |\det A_I|$\\
\midrule
2&12&13&3&1,949&12\\
2&48&53&7&515,367&61\\
3&12&13&5&344,908&102\\
3&32&37&11&405,205,698&1,066\\
3&64&67&16&65,061,419,952&7,680\\
4&12&13&6&43,843,180&1,936\\
4&30&31&13&232,929,760,400&18,500\\
5&17&17&9&167,433,728,000&23,661\\
5&23&23&12&34,155,205,413,840&224,106\\
\bottomrule
\end{tabular}
\end{center}
All nine gcds in this table are one. The exact-modularity tests use the following inputs; the actual largest minor is exactly the prescribed integer, and occurs once in each test. The finite tests use sharper direct norm bounds than the coarse uniform constant in Proposition~\ref{prop:exact}.
\begin{center}
\begin{tabular}{rrrr}
\toprule
$r$& prescribed $\Delta$ & number of columns & multiplicity of max\\
\midrule
2&100,003&101&1\\
3&1,000,001&31&1\\
4&1,000,000,000,039&31&1\\
5&1,000,000,000,000,037&17&1\\
\bottomrule
\end{tabular}
\end{center}
The program checks 162,973 maximal minors in each of its two passes, as well as all 17,367 integer points in seven test boxes (including very elongated ones) for the covering argument. The complete column lists, scalar multipliers, transformation matrices, and exact results are in \texttt{data/verification.json}. These finite checks supplement rather than replace the proofs.

The result resolves an exponent, not a leading coefficient. For $r=3$, the Gram minimum is of order $n^7$ and its squared normalized minors have mean of order at least $n^4$; the short-lift construction shows that this scale is achievable up to rank-dependent constants. For $r=4$ the corresponding exponents are $10$ and $6$. For rank two the conclusion $n^4$ agrees with the sharper existing asymptotic in ProveIt. By Hodge duality for primitive exterior products, as in the existing article, $E_r(n)=E_{n-r}(n)$ when both ranks are at least two. Thus the fixed-corank regime has the same exponent with $r$ replaced by that fixed corank. No uniform-in-rank estimate is asserted.

\section{Relationship to the GRH project and scope}
The upstream development at commit \texttt{72fead102f663d5c7a407964a13dc0cf9d977207} proves statements about Dirichlet $L$-functions in Lean, as described by its main file and the merged article \cite{upstream,proveit}. The proof architecture compares archimedean Gram data with integral structures and index losses. In particular the source theorem \texttt{New.TwistedLattice.covolume\_add\_index\_nonneg} explicitly combines a logarithmic Gram bound and an index correction with costs for completing a spanning family. The article's Section I abstracts the determinant part and explicitly raises the higher-rank problem solved here.

Our transfer consists of three steps: retaining the exact integer index normalization, finding a lower bound independent of the archimedean metric, and constructing a matching family using arithmetic reductions that guarantee independence. The homogeneous hyperplane cover is the additional rigidity needed when all maximal minors are nonzero. It does not occur as a claimed theorem of the upstream Lean development. All the arguments in this note take place in elementary linear algebra, lattice geometry, finite fields, and Bertrand's postulate. There is no dependency on GRH or any unproved analytic estimate.

\paragraph{Status and limits.} The merged article's proof claims are not independently audited in full here; the input theorem is not rebuilt. This note uses only the cited Gram question as motivation, and its proofs are self-contained apart from standard facts explicitly identified. It does not prove RH for further $L$-functions. The literature check identified the precise published question and later work on rank two and non-generic column numbers; no priority claim beyond this limited search is made. The leading constant for $r\ge3$, the refined behavior of minimizers, and the exact finite column numbers remain open here. The new arguments are human-readable and the computations are reproducible; they have not been formalized in Lean or independently peer reviewed.

\begin{thebibliography}{9}
\bibitem{proveit} ProveIt, \emph{Beyond the Riemann hypothesis for Dirichlet L-functions: consequences, transfers and limits}, repository article, Section I, \texttt{src:cor:gramrank} and \texttt{src:q:gram}. \href{https://github.com/VladimirReshetnikov/ProveIt/blob/main/Analysis/GeneralizedRiemannHypothesis/Article/sections/src.tex}{Source file} (consulted 27 November 2026; blob \texttt{b0ac2147291b7da3a74cdf6072cc557b1e1f5ee8}).
\bibitem{kks} B. Kriepke, G. M. Kyureghyan and M. Schymura, \emph{On the size of integer programs with bounded non-vanishing subdeterminants}, Journal of Combinatorial Theory, Series A \textbf{212} (2025), 106003. \href{https://doi.org/10.1016/j.jcta.2024.106003}{DOI}; \href{https://arxiv.org/abs/2309.03772}{arXiv:2309.03772}, especially Question 3.4.
\bibitem{ks} B. Kriepke and M. Schymura, \emph{On generic $\Delta$-modular integer matrices with two rows}, SIAM Journal on Discrete Mathematics (2026), \href{https://doi.org/10.1137/25M1745878}{DOI}; \href{https://arxiv.org/abs/2502.15394}{arXiv:2502.15394}. Related rank-two exact work, not a dependency of our argument.
\bibitem{upstream} L. Sernau and collaborators, \emph{grh-lean}, readable edition of the Lean proof of L. Alp\"oge with Claude, commit \texttt{72fead102f663d5c7a407964a13dc0cf9d977207}, \href{https://github.com/lsernau/grh-lean/tree/72fead102f663d5c7a407964a13dc0cf9d977207/grh}{repository snapshot}.
\end{thebibliography}
\end{document}
